\section[Fredholm determinants for q-Whittaker and shift-mixed periodic Schur measures]{Fredholm determinants for $\boldsymbol{q}$-Whittaker\\ and shift-mixed periodic Schur measures}
\label{s2}

The main purpose of this section is to give Fredholm determinant formulas for both sides of~\eqref{t12-3}; see Propositions~\ref{p22} and~\ref{p25}.
In Section~\ref{s3}, we will show that the two Fredholm determinants are in fact the same,
as stated in Theorem~\ref{t31}. In Section~\ref{s23}, we also give a proof of equivalence between Theorems~\ref{t1} and~\ref{t12}.

Hereafter, we consider set of variables $a_1,\dots,a_N$, $b_1,\dots,b_N$, i.e., we set $M=N$ in the notation used in Section~\ref{s1}. The general case $M \neq N$ of each result is recovered by considerations similar to the one reported in Remark~\ref{rem:intro}. We will denote the maximum (resp.\ minimum) value of $a_i$ and $b_i$ $i=1,\dots,N$ by $a_{\max}$, $b_{\max}$ (resp.\ $a_{\min}$, $b_{\min}$). We also always omit the factor $1/\big(2\pi \sqrt{-1}\big)$ in contour integrals.

\subsection[Fredholm determinant formula for the q-Whittaker measure]{Fredholm determinant formula for the $\boldsymbol{q}$-Whittaker measure}
\label{s21}

We focus on the expectation value $\E[1/(\zeta q^{-\mu_1};q)_{\i}]$ with respect to the $q$-Whittaker measure defined by \eqref{s13}. This quantity is the $q$-Laplace transform of the marginal $\mu_1$ and it admits explicit Fredholm determinant expressions \cite{BorodinCorwin2014Mac,BoCoFeVe2015,BufetovMucciconiPetrov2019,ImamuraSasamoto2019qTASEPtheta,MatveevPetrov2015}. For our arguments, we will make use of the formula that first appeared in \cite{ImamuraSasamoto2019qTASEPtheta}, which is conceptually different from the ones of the other references.

The following proposition is obtained by applying a few results
in \cite{ImamuraSasamoto2019qTASEPtheta}.


\begin{Proposition}\label{p22}
Let $\zeta\in\C\setminus\{q^n\}_{n\in\Z}$, and
$a_i, b_i\in\R_{>0}$, $i=1,\dots,N$ satisfy $a_{\max}b_{\max}<1$ and~${qa_{\max}<a_{\min}}$.
Let $\mu_1$ be the first row of $\mu$ distributed according to the $q$-Whittaker measure~\eqref{s13}.
We have
\begin{align}
\label{p221}
\E\left[\frac{1}{(\zeta q^{-\mu_1};q)_{\i}}\right]
=
\det(1-f_{\zeta} K)_{\ell^2(\Z)},
\end{align}
where
$f_{\zeta}(m)=\frac{-\zeta q^m}{1-\zeta q^m}$
and
\begin{align}
\label{p222}
&K(m_1,m_2)=
\int_C\frac{{\rm d}z}{z}\int_{D}\frac{{\rm d}w}{w}
g_{a,b}(z,w;m_1,m_2),
\\
&
g_{a,b}(z,w;m_1,m_2)
=
\frac{w^{m_2}}{z^{m_1}}
\frac{w}{z-w}
\prod_{i=1}^N
\frac{(a_i z;q)_{\i}}{(a_iw;q)_{\i}}
%\cdot
%\prod_{j=1}^N
\frac{(b_i/w;q)_{\i}}{(b_i/z;q)_{\i}},
\label{p223}
\end{align}
where the contour $C$ encloses anticlockwise
the poles at $z=0, b_iq^j$, $i=1,\dots,N$,
$j\in\Z_{\ge 0}$
whereas~$D$ encloses anticlockwise
the poles at
$w=1/a_i$, $i=1,\dots,N$
but not the ones at $w=0, 1/a_iq^j$,
$i=1,\dots,N$, $j\in\Z_{\ge 1}$.
They are illustrated in Figure~{\rm\ref{f1}}.
\end{Proposition}

\begin{figure}[ht]
 \centering
 \includegraphics[scale=1]{fig1.pdf}
 %\includegraphics[]{contours_IMS21_1.pdf}
 \caption{The contours $C$, $D$, and $\bar{D}$ appearing in \eqref{p222}
 and \eqref{t252}. ``$\bullet$''s and ``\boldcross''s represent the poles of
 $w$ and $z$ respectively.}
 \label{f1}
\end{figure}


To prove Proposition~\ref{p22}, we recall the discussions in \cite[Section 3]{ImamuraSasamoto2019qTASEPtheta}. There we consider a~two-sided $q$-Whittaker measure, which is a~probability measure on the set of signatures,
\begin{align*}
 \S_N=\big\{\l=(\l_1,\dots,\l_N)\in \Z^N\mid \l_1\ge\cdots\ge \l_N\big\}.
\end{align*}
Note that $\l_i$, $i=1,\dots,N$ can be negative. Such probability measures on the extended space can be realized by replacing the function $Q_{\mu}(x)$ in \eqref{s13} with $Q_{\mu}(\rho)$, where $\rho$ is a certain specialization of the algebra of symmetric functions.
Here we state only a particular case which corresponds to \cite[Definition 3.4]{ImamuraSasamoto2019qTASEPtheta}, setting $a_i=x_i$, $i=1,\dots,N$, $t=0$ and $\a_j=b_j$, $j=1,\dots,N$. With this choice $M_{q\text{W}}(\mu)$ in \eqref{s13} becomes
\begin{align}
\label{p224}
P_{\l}(x_1,\dots,x_N)Q_{\l}(\rho^{-}_b)/\Pi^{(-)}(x;b),
\end{align}
where $\l\in\S_N$, the parameters $x_i$, $b_i$,
satisfy $b_{\max}/x_{\min}<1$ and
\begin{align}
Q_\l(\rho^-_b)
 =\prod_{i=1}^{N-1}\big(q^{\l_i-\l_{i+1}+1};q\big)_{\infty}
\int_{C^N}\prod_{i=1}^N\frac{{\rm d}z_i}{z_i}\cdot
P_\l (1/z)
 \Pi^{(-)}(z;b)m_N^q(z).
\label{defQ}
\end{align}
Here $C$ is the same contour as the
one in \eqref{p222}, i.e., it encloses
positively
the poles at $z=0, b_iq^j$, $i=1,\dots,N$
and $j\in\Z_{\ge 0}$, $1/z$ in $P_\l$ is a shorthand notation for
$(1/z_1,\dots,1/z_N)$ and
\begin{align*}%\label{Sk}\label{d14}
&m_N^q(z)=\frac{1}{N!}\prod_{1\le i<j\le N}(z_i/z_j;q)_{\infty}(z_j/z_i;q)_{\infty},
\qquad
\Pi^{(-)}(z;b)=\prod_{i,j=1}^N\frac{1}{(b_i/z_j;q)_{\infty}}.
\end{align*}
For more detail see \cite[Appendix B]{ImamuraSasamoto2019qTASEPtheta}.

From \cite[Lemma 3.3, equation~(3.15)]{ImamuraSasamoto2019qTASEPtheta}, one sees that $Q_{\l}(\rho_b^-)=0$ unless
$0\ge \l_1\ge\cdots \ge \l_N$ and thus the measure
\eqref{p224} has support only on the negative region. Although the measures~\eqref{s13} and \eqref{p224} have different support, they enjoy the following relation:
\begin{Lemma} \label{l24}
Set $a_i=1/x_i$, $i=1,\dots,N$ and define the signature $\mu$ by the negation of $\l$,
$\mu:=(-\l_N,-\l_{N-1},\dots,-\l_1)$.
 Then we have
\begin{align}
\label{2-12d}
 P_{\l}(x_1,\dots,x_N)Q_{\l}(\rho^{-}_b)/\Pi^{(-)}(x;b)
=
M_{q{\rm W}}(\mu),
 \end{align}
 and in particular
\begin{align}
\label{2-12}
 \E_{q{\rm W}(x;\rho_b^-)}\left[\frac{1}{\big(\zeta q^{\l_N};q\big)_{\i}}\right]
 =
 \E_{q{\rm W}(a;b)}\left[\frac{1}{(\zeta q^{-\mu_1};q)_{\i}}\right],
\end{align}
 where \smash{$\E_{q{\rm W}(x;\rho_b^-)}[\cdot]$}
 and $\E_{q{\rm W}(a;b)}[\cdot]$
 are averages over the measures
 \eqref{p224} and \eqref{s13}, respectively.
\end{Lemma}
\begin{proof}
 Equation \eqref{2-12d} follows straightforwardly from a combination of definition \eqref{p224}
 and relations \cite[equations~(3.14) and (3.15)]{ImamuraSasamoto2019qTASEPtheta},
 \[%\label{2-14}
 P_{\l}(x_1,\dots,x_N)=P_{\mu} (a_1,\dots,a_N),
 \qquad
 Q_{\l}(\rho^-_b)=Q_{\mu}(b_1,\dots,b_N),
 \]
 once we recall hypothesis relating $a_i$ (resp.\ $\mu$)
 to $x_i$ (resp.\ $\l$).
\end{proof}

\begin{proof}[Proof of Proposition~\ref{p22}]
 It follows immediately from Lemma~\ref{l24}
 and the Fredholm determinant formula
 for the left-hand side of \eqref{2-12}.
 Changing the notations $a_i\rightarrow x_i$,
 $\alpha_j\rightarrow b_j$,
 $v\rightarrow w$,
 $n\rightarrow -m_1$, $m\rightarrow -m_2$
 and setting $N=M$ and $t=0$
 in~\cite[equation~(4.37)]{ImamuraSasamoto2019qTASEPtheta},
 we obtain the Fredholm determinant formula
\begin{align}
\label{p225}
\E_{q\text{W}(x;\rho_b^-)}\left[\frac{1}{\big(\zeta q^{\l_N};q\big)_{\i}}\right]
=
\det(1-f_{\zeta} M)_{\ell^2(\Z)}.
\end{align}
Here $f_{\zeta}(m)$ is defined above~\eqref{p222}
and
\begin{align*}%\label{t133}
M(m_1,m_2)=
\int_C\frac{{\rm d}z}{z}\int_{D}\frac{{\rm d}w}{w}
\frac{w^{m_2}}{z^{m_1}}
\frac{w}{z-w}
\prod_{i=1}^N
\frac{(z/x_i;q)_{\i}}{(w/x_i;q)_{\i}}
\frac{(b_i/w;q)_{\i}}{(b_i/z;q)_{\i}},
\end{align*}
where $C$ is the same contour as \eqref{defQ} and $D$ encloses positively
$x_i,~i=1,\dots,N$.
Combining~\eqref{2-12} with~\eqref{p225},
we get~\eqref{p221}.
\end{proof}

\subsection[Fredholm determinant formula for the shift-mixed periodic Schur measure]{Fredholm determinant formula \\
for the shift-mixed periodic Schur measure}\label{s22}

We have discussed in Section~\ref{s1} how the periodic Schur measure $M_{\text{pS}}(\l)$ \eqref{1-8} becomes a determinantal point process once all the rows of $\lambda$ get shifted by an independent random quantity $S$ distributed as \eqref{1-11}. This was originally observed by Borodin in \cite{Borodin2007PeriodicSchur} and the measure describing the process
\begin{align}
\label{216}
(\l_1+S,\l_2+S,\dots)
\end{align}
takes the name of \emph{shift-mixed periodic Schur measure}. The probability distribution of the first element $\l_1+S$ can be represented as a single Fredholm determinant and this is reported in the next proposition.


\iffalse
[Old: Here we give a Fredholm determinant formula for a deformed version of the periodic Schur measure $M_{\text{pS}}(\l)$\eqref{1-8} called the shift-mixed periodic Schur measure. The point process $\l$ followed by
the marginal of $M_{\text{pS}}(\l)$ is not a determinantal point process (DPP). In~\cite{Borodin2007PeriodicSchur}, Borodin found out that this situation can be modified by considering a random shift, i.e., the point process shifted by the integer-valued random variable $S$,
\begin{align}
\label{216}
(\l_1+S,\l_2+S,\dots)
\end{align}
becomes the DPP. Here $S$ is independent of $M_{\text{pS}}$ and follows the distribution \eqref{1-11}. The product measure of \eqref{1-8} and \eqref{1-11} is called the {\it shift-mixed periodic Schur measure}. As mentioned above, this is a determinantal point process. Accordingly probability distribution of the first element with shift, $\l_1+S$, can be represented as a single Fredholm determinant. The following proposition is a direct consequence of this fact.]
\fi


\begin{Proposition}\label{p25}
 Let $\l_1+S$ be the first element of \eqref{216} distributed according to the shift-mixed periodic Schur measure.
 For $k\in\Z_{\ge 0}$, we have
\begin{align} \label{p251}
 \P (\l_1+S\le k )
 =\det (1+f L )_{\ell^2(\Z)},
\end{align}
 where
 $f(m)$ is given in terms of
 $f_{\zeta}(m)$ defined above \eqref{p222},
\begin{align}
\label{p251-1}
 f(m)=f_{\zeta}(m)|_{\zeta=-tq^{1/2+k}}
 =\frac{tq^{1/2+k+m}}{1+tq^{1/2+k+m}}
\end{align}
and
by using $g_{a,b}(z,w,m_1,m_2)$~\eqref{p223},
the kernel
$L(m_1,m_2)$ is defined by
\begin{align}
&L(m_1,m_2)=
\int_C\frac{{\rm d}z}{z}\int_{\bar{D}}\frac{{\rm d}w}{w}
g_{a,b}(z,w,m_1,m_2).
\label{t252}
\end{align}
Here contour $C$ is oriented counterclockwise and it encloses the origin with its radius
$r$ satisfying~${b_{\max}<r<1/a_{\max}}$
whereas $\bar{D}$ encloses $C$
anticlockwise.
$C$ and $\bar{D}$ are depicted in Figure~{\rm\ref{f1}}
given below Proposition~{\rm\ref{p22}}.

\end{Proposition}


\begin{proof}
From \cite[Theorem 2.2]{Borodin2007PeriodicSchur}, we have, for $k\in\Z_{\ge 0}$,
\begin{align*}%\label{111}
\P(\l_1+S\le k)
=\det\big(1-\tilde{L}\big)_{\ell^2(k+1,k+2,\dots)},
\end{align*}
where the kernel $\tilde{L}$ is defined as
\begin{align}
\tilde{L}(\ell_1,\ell_2)={}&\int_{|w|=r}\frac{{\rm d}w}{w}\nonumber\\
&\times\int_{|z|=r'}\frac{{\rm d}z}{z}\frac{z^{\ell_2-1}}{w^{\ell_1-1}}
\prod_{i=1}^N
\frac{(a_iz;q)_\i}{(a_i w;q)_{\i}}
\frac{(b_i/w;q)_\i}{(b_i/z;q)_\i}
\cdot
\sum_{m\in \Z}
\frac{t q^{\frac12+m}}{1+t q^{\frac12+m}}
\left(\frac{w}{z}\right)^m,
\label{112}
\end{align}
and $r$, $r'$ satisfy\footnote{In~\cite{BeteaBouttier2019PeriodicSchur,Borodin2007PeriodicSchur},
the conditions for $r$, $r'$ are
given by using the analyticity
of the integrand for the periodic
Schur process with positive specializations.
This condition can be translated into
the explicit one \eqref{t111}
in our case~\eqref{1-8}.
}
\begin{align}
b_{\max}<r,r'<\frac{1}{a_{\max}},\qquad 1<\frac{r}{r'}<q^{-1}.
\label{t111}
\end{align}
We will now show that $\det\big(1-\tilde{L}\big)_{\ell^2(k+1,k+2,\dots)} =\det(1+fL)_{\ell^2(\Z)}$. Note that the kernel $\tilde{L}(\ell_1,\ell_2)$ \eqref{112} can be written as the kernels of the operators $A$ and $B$, which map between $(k+1,k+2,\dots)$ and $\mathbb{Z}$:
\begin{align*}%\label{t133-1}
\tilde{L}(\ell_1,\ell_2)=\sum_{m\in\Z}A(\ell_1,m)
B(m, \ell_2),
\end{align*}
where
\begin{align*}%\label{t134}\label{t135}
A(\ell,m)&=\int_{|w|=r}\frac{{\rm d}w}{w} w^{m-\ell+1}
\prod_{i=1}^N
\frac{(b_i/w;q)_\i}{(a_i w;q)_{\i}},
\\
B(m,\ell)&=
\frac{t q^{1/2+m}}{1+t q^{1/2+m}}
\int_{|z|=r'}\frac{{\rm d}z}{z} z^{\ell-1-m}
\prod_{i=1}^N
\frac{(a_iz;q)_\i}{(b_i/z;q)_\i}.
\end{align*}
By a basic property of the determinant, we have
\begin{align*}%\label{t135-1}
\det(1-AB)_{\ell^2(k+1,k+2,\dots)}=\det(1-BA)_{\ell^2(\Z)},
\end{align*}
and evaluating the kernel in the right-hand side, we find
\begin{gather*}%\label{t136}
(BA)(n_1,n_2)
\\
\qquad=
\frac{t q^{n_1+\frac12}}{1+t q^{n_1+\frac12}}
\int_{|z|=r'}\frac{{\rm d}z}{z}
\int_{|w|=r}\frac{{\rm d}w}{w}
\frac{w^{n_2}}{z^{n_1}}
\prod_{i=1}^N
\frac{(a_iz;q)_\i}{(a_i w;q)_{\i}}
\frac{(b_i/w;q)_\i}{(b_i/z;q)_{\i}}
\cdot
\sum_{\ell=k+1}^{\i}\left(\frac{z}{w}\right)^{\ell-1}
\\
\qquad=
\frac{-t q^{\frac12+k+m_1}}{1+t q^{\frac12+k+m_1}}
\int_{|z|=r'}\frac{{\rm d}z}{z}
\int_{|w|=r}\frac{{\rm d}w}{w}
\frac{w^{m_2}}{z^{m_1}}
\prod_{i=1}^N
\frac{(a_iz;q)_\i}{(a_i w;q)_{\i}}
\frac{(b_i/w;q)_\i}{(b_i/z;q)_{\i}}
\cdot
\frac{w}{z-w}
\\
\qquad=
-f(m_1)L(m_1,m_2).
\end{gather*}
Here in the second equality we set $m_i=n_i-k$ for $i=1,2$.
\end{proof}



\subsection[Equivalence of Theorems 1.1 and 1.3]{Equivalence of Theorems \ref{t1} and \ref{t12}}
\label{s23}



This subsection is devoted to proving that Theorem~\ref{t1} is equivalent to Theorem~\ref{t12}. Since equivalence between Theorems~\ref{t1} and~\ref{t12}\,(a) is straightforward, the only nontrivial part is to show the equivalence between Theorem~\ref{t12}\,(a) and (b), i.e., \eqref{t12-1} and \eqref{t12-3}. This is a~consequence of the computation reported in the following lemma.
\begin{Lemma}
\label{l27}
For $n\in\Z$, $0<q<1$ and
$t>0$, we have
 \begin{align}
\label{2-31}
 \P(\chi+S\le n)
 =
 \frac{1}{\bigl(-t q^{\frac12+n};q\bigr)_{\i}},
\end{align}
 where $\chi$ and $S$ are independent random variables with distributions
 \eqref{t11-3} and \eqref{1-11}, respectively.
\end{Lemma}
 \begin{proof}
 We write the left-hand side of \eqref{2-31} as
 \begin{align}
 \label{2-34}
 \P(\chi+S\le n)
 =
 \sum_{\ell\in\Z} \P(S=\ell)\P(\chi\le n-\ell).
 \end{align}
 Applying \eqref{1-11} and
 \begin{align*}
 \P(\chi\le m)=
 \begin{cases}
 (q;q)_{\i}/(q;q)_m,& m\in\Z_{\ge 0},
 \\
 0,& m\in\Z_{< 0},
 \end{cases}
 \end{align*}
 which follows from \eqref{t11-3} and the fact $\sum_{n=0}^mq^n/(q;q)_n=1/(q;q)_m$, $m=0,1,2,\dots$, which can be checked easily by mathematical induction, to the right-hand side of \eqref{2-34}, we find
 \begin{align*}
 \P(\chi+S\le n)
 =
 \frac{1}{\theta\bigl(-q^{1/2}t\bigr)}\sum_{\ell=-\infty}^n
 \frac{t^{\ell}q^{\frac{\ell^2}{2}}}
 {(q;q)_{n-\ell}}
 =
 \frac{t^nq^{\frac{n^2}{2}}\bigl(-q^{\frac12-n}t^{-1};q\bigr)_{\i}}{\theta\bigl(-q^{1/2}t\bigr)},
 \end{align*}
 where, in the second equality we used a version of the $q$-binomial theorem \eqref{a2-4}.
 Noting from the definitions of the $q$-Pochhammer symbol \eqref{a2-1} and \eqref{a2-1d}, the numerator of the last expression becomes
 \begin{align*}
 t^n q^{n^2/2} \bigl(-q^{\frac12-n}t^{-1};q\bigr)_\i
=\bigl(-t q^{1/2};q\bigr)_n \bigl(-q^{1/2}/t;q\bigr)_{\i},
 \end{align*}
 we arrive at the right-hand side of \eqref{2-31}.
 \end{proof}

 By using \eqref{2-31}, we see that
 the left-hand side of \eqref{t12-3}
 is expressed as
 \begin{align*}
 \E
 \left[ \frac{1}{\bigl(-t q^{\frac12+n-\mu_1}\bigr)}\right]
 =\P(\mu_1+\chi+S\le n).
 \end{align*}
 Thus \eqref{t12-3} can be rewritten as
 \begin{align*}% \label{2-33}
 \P(\mu_1+\chi+S\le n)
 =
 \P(\l_1+S\le n),
 \end{align*}
 which is clearly equivalent to \eqref{t12-1}.

\section[Equivalence of two Fredholm determinants by shift of contours]{Equivalence of two Fredholm determinants\\ by shift of contours} \label{s3}

In this section, we will force some conditions on parameters $a_i$, $b_i$,
$i=1,\dots,N$. Namely we assume that $a_i,b_i\in (0;1)$, $i=1,\dots,N $ are distinct and ordered as $a_1>a_2>\cdots > a_N$ (hence~${a_{\max}=a_1}$ and $a_{\min}=a_N$)
and moreover
\[
 a_{1}/a_{N}<q^{-\frac12+\e},
\]
for a fixed $\e \in (0,1/2)$. Recall that for the sake of our main result Theorem~\ref{t1} these restrictions do not constitute a problem, see Remark~\ref{rem:intro}.


\subsection{Restating the main theorem as a determinant identity}
\label{s31}
In the previous section, precisely in Propositions~\ref{p22} and~\ref{p25}, we have seen that both sides of~\eqref{t12-3} can be written as Fredholm determinants. By this equality \eqref{t12-3} is equivalent to the following identity for the two Fredholm determinants.
\begin{Theorem}
\label{t31}
Let $K(m_1,m_2)$ be the kernel given in ~\eqref{p222} and $L(m_1,m_2)$ in ~\eqref{t252}.
Also let~$f$ be given by \eqref{p251-1}.
Then we have
\begin{align}
\label{31}
\det(1-f K)_{\ell^2(\Z)}
=
\det(1+fL)_{\ell^2(\Z)}.
\end{align}
\end{Theorem}
In this section, we will elaborate the proof of Theorem~\ref{t31}.
Comparing the expressions of the two kernels in~\eqref{p222} and in~\eqref{t252}, we notice that both of them are written in the double integral forms with the common integrand $g_{a,b}(z,w;m_1,m_2)$, given by \eqref{p223}. The only difference between them is the integration contour of $w$: the contour $D$ in $K(m_1,m_2)$ encloses $1/a_j$, $j=1,\dots,N$, while $\bar{D}$ in $L(m_1,m_2)$ encloses the origin and the contour of $z$. See Figure~\ref{f1}.


Result of Theorem~\ref{t31} is not trivial outside of the case $q=0$, when the kernel themselves are equal, i.e., $K(m_1,m_2)=-L(m_1,m_2)$ for all $m_1,m_2\in\Z_{\le 0}$. This can be shown as follows. When $q=0$, we notice that $f(m)$ becomes the indicator function of $\mathbb{Z}_{\le 0}$ and the integrand~${g_{a,b}(z,w;m_1,m_2)}$ in \eqref{p223} becomes
\begin{align*}
 g_{a,b}(z,w;m_1,m_2)
 =\frac{w^{m_2}}{z^{m_1}}
\frac{w}{z-w}
\prod_{i=1}^N
\frac{(1-a_i z)}{(1-a_iw)}
\frac{(1-b_i/w)}{(1-b_i/z)},
\end{align*}
and we see that the poles in the $w$ variable include only $w=0, z, 1/a_i$, $i=1,\dots,N$ and in particular $\infty$ is not a pole since $m_2 \le 0$. Thus using Cauchy's integral theorem, we see that the contour $D$ can be deformed to $\bar{D}$ without changing the value of the kernel up to sign.

For general $0<q<1$, this is no longer the case. In this case, equality does not hold for the kernels, but it only does at the level of the Fredholm determinants \eqref{31}. To prove this we consider a sequence of changes of the contours of $w$ starting from $D$, and including, sequentially, poles of the form $1/a_j q^k$, $j=1,\dots,N$ starting from $k=0,1,\dots$ and finally including the singularity at infinity. We find that such modifications of the kernel do not affect its Fredholm determinant, so that, once all the singularities have been added in the contour one can switch such contour to $\bar{D}$ by using the Cauchy theorem.


Notations for the series of contours just described, their corresponding kernels and Fredholm determinants are provided next.
\begin{Definition}\label{d31}
For $\ell\in\mathbb{Z}_{\ge 0}$, we define the kernels
\begin{align}
\label{d31-3}
&K_\ell (m_1,m_2)=
\int_C\frac{{\rm d}z}{z}\int_{D_\ell}\frac{{\rm d}w}{w} g_{a,b}(z,w;m_1,m_2),
\\
\label{d31-4}
&K_{\i}(m_1,m_2)
=
\int_C\frac{{\rm d}z}{z}\int_{D_{\i}}\frac{{\rm d}w}{w} g_{a,b}(z,w;m_1,m_2).
\end{align}
The function $g_{a,b}$ is given by \eqref{p223} and the contour $C$ encloses the origin and it has radius larger than $b_{\max}$ as in \eqref{p222}, \eqref{t252}. The integration contours $D_\ell$ and $D_{\infty}$ are illustrated in Figure~\ref{fig2}.
Note $D_\ell$ encircles poles $1/q^{i}a_j$, $i=0,\dots,\ell$, $j=1,\dots,N$ and no other singularities lie inside~$D_{\ell}$.
Recalling the function $f(m)$ defined by \eqref{p251-1}, we further define the Fredholm determinants
\begin{align}
\label{d31-1}
&F_\ell:=\det(1-fK_\ell)_{\ell^2(\Z)}\qquad \text{for}\quad \ell=0,1,2,\dots,
\\
\label{d31-2}
&F_{\infty}:=\det(1-f K_{\i})_{\ell^2(\Z)}.
\end{align}
\end{Definition}

\begin{figure}[ht]
 \centering
 \includegraphics[scale=1]{fig2.pdf}
 \caption{Contours $D_{\ell}$, $\tilde{D}_\ell$, and $D_{\i}$
 are illustrated.
 We set $c_{\ell}$
 to be the center of the interval $\big[\frac{1}{a_{N}q^\ell},
 \frac{1}{a_{1} q^{\ell+1}}\big]$, i.e., $c_{\ell}=\frac{1}{2q^\ell}\big(\frac{1}{a_{1}q}+\frac{1}{a_{N}}
 \big)$.}
 \label{fig2}
\end{figure}


Note that by definition, we see $K_0(m_1,m_2)=K(m_1,m_2)$. Thus by Proposition~\ref{p222}, we immediately see $F_0=\det(1-fK)_{\ell^2(\Z)}$. In \cite[Proposition 4.7]{ImamuraSasamoto2019qTASEPtheta}, we showed that the kernel~${K(m_1,m_2)}$ is trace-class, thus $\det(1-fK)_{\ell^2(\Z)}$ is well defined.
This fact will be generalized to $F_{\ell}$ for any~${\ell\in\Z_{\ge 0}}$ and also $F_{\i}$ in Lemma~\ref{l32}\,(iii).

For this purpose and also for later use,
we will give different representations for
the kernels~\eqref{d31-3} and \eqref{d31-4}.
 To simplify notation, we introduce a relabelling of all the inverse of poles~$a_i q^j$ in
decreasing order and call them $\tilde{a}_r,r\in\Z_{>0}$. Namely, we write
\begin{align}
\label{s32-1}
 \tilde{a}_{r}=a_{k}q^u,
\end{align}
where $r=uN+k$ with $r\in\Z_{>0}$, $k\in\{1,\dots,N\}$ and $u\in\Z_{\ge 0}$.
Since we assumed $a_1>a_2>\cdots >a_N$ without loss of generality, $\tilde{a}_1>\tilde{a}_2>\cdots$ holds.

For $m\in\Z$ and $r\in\Z_{\ge 0}$,
we define two infinite-dimensional matrices,
$A(m;r)$, $B(r;m)$, as
\begin{align}
 \label{3-15d}
 A(m;r)&=f(m)\int_C\frac{{\rm d}z}{z^{1+m}}
 \frac{1}{z-\tilde{a}_{r}^{-1}}
 \prod_{i=1}^N\frac{(a_i z;q)_{\i}}{(b_i/z;q)_{\i}},
 \\
 \label{3-16d}
 B(r;m)&=\tilde{a}_{r}^{-m}
 \underset{w=\tilde{a}_{r}^{-1}}{\text{Res}}
 \prod_{i=1}^N\frac{(b_i/w;q)_{\i}}{(a_i w;q)_{\i}}.
\end{align}
For $\ell\in\Z_{\ge 0}$, we also introduce their truncated version, $A_{\ell}(m;r)$, $B_{\ell}(r;m)$, defined by
\[%\label{3-8}
 A_{\ell}(m;r):=A(m;r) \mathbf{1}_{\{1,\dots, N(\ell+1)\}}(r),\qquad
 B_{\ell}(r;m):= \mathbf{1}_{\{1,\dots, N(\ell+1)\}}(r)B(r;m),
\]
where for a set $A\subset\Z_{\ge0}$, $\mathbf{1}_A$ is the indicator function of $A$. By evaluating the poles inside $D_\ell$ of~\eqref{d31-3}, we easily find
\begin{align}
\label{3-10}
 &f(m_1)K_{\ell}(m_1,m_2)=\sum_{r=1}^{\infty}A_{\ell}(m_1;r)B_{\ell}(r;m_2).
\end{align}
Below in the proof of Proposition~\ref{p33}\,(iii), we will also see that
\begin{align}
\label{3-11}
 f(m_1)K_{\infty}(m_1,m_2)=\sum_{r=1}^{\i}A(m_1;r)B(r;m_2).
\end{align}
In the following discussions, it is also useful to consider the conjugated matrices
\begin{align}
 \label{3-16dd}
 \tilde{A}(m;r)&=\t(m) A(m;r)\sigma(r),
 \qquad
 \tilde{B}(r;m)=\sigma^{-1}(r) B(r;m)\t^{-1}(m),
\end{align}
where
\begin{align}
\label{3-12d}
 \t(m)=
 \begin{cases}
 a_{1}^{-m} q^{-\frac{m^2}{2N}-\frac{m}{2}+\e m},
 & m\ge 0,\\
 1,& m<0,
 \end{cases}
 \qquad
 \sigma(r)= q^{-(1-\omega)u}
\end{align}
for $\e\in(0,1/2)$ and $\o\in(0,1/2-\e)$, and also
their truncated versions $\tilde{A}_{\ell}(m;r)$, and
$\tilde{B}_{\ell}(r;m)$,
\begin{align}
\label{3-15at}
 \tilde{A}_{\ell}(m;r):=\tilde{A}(m;r) \mathbf{1}_{\{1,\dots, N(\ell+1)\}}(r),\qquad
 \tilde{B}_{\ell}(r;m):= \mathbf{1}_{\{1,\dots, N(\ell+1)\}}(r)\tilde{B}(r;m).
\end{align}
Note that the value of the
Fredholm determinant would not change when we replace
$A_\ell$, $B_\ell$,~$A$,~$B$ by their conjugated versions,
$\tilde{A}_\ell$, $\tilde{B}_\ell$, $\tilde{A}$, $\tilde{B}$
in the kernel \eqref{3-10}, \eqref{3-11}.

\begin{Proposition}\label{p33}
Assume $a_i,b_i\in(0,1)$, $i=1,\dots,N$, $a_1>\cdots>a_N$, and $a_{1}/a_{N}<q^{-\frac12+\e}$ with $\e\in(0,1/2)$.
\begin{itemize}\itemsep=0pt
 \item[{\rm (i)}]
The operators with
the kernels, $\tilde{A}(m;r)$, $\tilde{B}(r;m)$,
$\tilde{A}_\ell(m;r)$, $\tilde{B}_\ell(r;m)$
for $m\in\Z$ and $r,\ell\in\Z_{\ge 0}$
are of Hilbert--Schmidt class.
 \item[{\rm (ii)}] $K_{\i}(m_1,m_2)$
 can be expressed as the limit $\ell\rightarrow\i$ of $K_{\ell}(m_1,m_2)$,
 \begin{equation}
 \label{p33-1}
 K_{\i}(m_1,m_2)=\lim_{\ell\rightarrow\i}K_\ell(m_1,m_2).
 \end{equation}
 \item[{\rm (iii)}]
 The Fredholm determinants $F_\ell $ and $F_{\infty}$, defined by~\eqref{d31-1} and~\eqref{d31-2}, respectively, are well defined.
\end{itemize}
\end{Proposition}
\begin{proof}
 (i) By the estimates of $\tilde{A}(m;r)$, $\tilde{B}(r;m)$ in Corollary~\ref{cb2},
 it is easy to check
 the conditions of the Hilbert--Schmidt operators
\[
 \sum_{m\in\Z}\sum_{r=1}^{\i}\big|\tilde{A}(m;r)\big|^2<\i,
 \qquad
 \sum_{m\in\Z}\sum_{r=1}^{\i}\big|\tilde{B}(r;m)\big|^2<\i.
\]
Obviously, the same relations also hold for $\tilde{A}_{\ell}$, $\tilde{B}_{\ell}$.

 For (ii), let us introduce the notation,
\begin{align}
\label{3-26}
 K(m_1,m_2;\sharp)=
 \int_C\frac{{\rm d}z}{z}
 \int_{\sharp}\frac{{\rm d}w}{w}
 g_{a,b}(z,w,m_1,m_2), \qquad m_1,m_2\in\Z,
\end{align}
 where $g_{a,b}(z,w,m_1,m_2)$ is defined by \eqref{p223} and $\sharp\in\big\{D_{\ell}, \tilde{D}_{\ell}\big\}_{\ell=0,1,\dots}\cup\{D_{\i}\}$ represents a contour for $w$. They are illustrated in Figure~\ref{fig2}. By definition, one sees
\[%\label{s34-1}
 K_\ell(m_1,m_2)=K(m_1,m_2;D_\ell),\qquad
 K_{\i}(m_1,m_2)= K(m_1,m_2;D_\i),
\]
Using~\eqref{3-26}, we show \eqref{p33-1}.
 Noting $D_{\i}=D_\ell\cup\tilde{D}_\ell$, we see that
\begin{align*}%\label{s34-3}
 K(m_1,m_2;D_{\i})
 &=
 K(m_1,m_2;D_\ell)+ K\big(m_1,m_2;\tilde{D}_\ell\big)
 \\
 &=
 \lim_{\ell\rightarrow\i}K(m_1,m_2;D_\ell)
 +
 \lim_{\ell\rightarrow\i}K\big(m_1,m_2;\tilde{D}_\ell\big)
 =\lim_{\ell\rightarrow\i}K(m_1,m_2;D_\ell).
\end{align*}
 Here in the second equality, we used the
 fact that $D_\ell\cup\tilde{D}_\ell$ $(=D_{\i})$
 does not depend on $\ell$ while in the third equality we used the
 fact \smash{$\lim_{\ell\rightarrow\i}K\big(m_1,m_2;\tilde{D}_\ell\big)=0$}. This follows from estimates of Lemma \ref{lc1}\,(iii): the contour $\tilde{D}_\ell$ can be divided into three
 parts \smash{$\tilde{D}_\ell=\tilde{D}_\ell^{(1)}\cup \tilde{D}_\ell^{(2)}\cup \tilde{D}_\ell^{(3)}$}, where~\smash{$\tilde{D}_\ell^{(1)}$} \big(resp.\ \smash{$\tilde{D}_\ell^{(3)}$}\big) are horizontal half line from~${+\infty+ d\sqrt{-1}}$ to $c_{\ell}+d\sqrt{-1}$ (resp.\ from $c_{\ell}-d\sqrt{-1}$ to $+\infty-d\sqrt{-1}$) while \smash{$\tilde{D}_\ell^{(2)}$} is vertical line segment from $c_{\ell}+d\sqrt{-1}$ to $c_{\ell}-d\sqrt{-1}$. We get $\lim_{\ell\rightarrow\i}K\big(m_1,m_2;\tilde{D}_\ell\big)=0$ applying \eqref{lc1-3} for the cases \smash{$\tilde{D}_{\ell}^{(1)}$} and \smash{$\tilde{D}_{\ell}^{(3)}$} and \eqref{lc1-3'} for the case \smash{$\tilde{D}_{\ell}^{(2)}$} to the part \smash{$\prod_{i=1}^N(a_iw;q)_{\infty}^{-1}$} in the integrand \eqref{p223}.

(iii)
 By equations~\eqref{3-10}, \eqref{3-15at} and (i), we see that the operators with the kernel
 \[\t(m_1)f(m_1) K_\ell(m_1,m_2)\t^{-1}(m_2),\]
 $\ell=0,1,2,\dots$ are trace-class. Since the factor $\t(m_1)/\t(m_2)$ does not affect the value of the determinant, then $F_\ell$ in \eqref{d31-1} is well defined.
 Next combining \eqref{3-10} and~\eqref{p33-1}, we have \eqref{3-11}. Then by the same arguments used above, the operator with the kernel $\t(m_1)K_\i (m_1,m_2)\t^{-1}(m_2)$ is trace-class and $F_\infty$ in \eqref{d31-2} is well defined.
\end{proof}

Theorem~\ref{t31} can be immediately proved by the following lemma.

\begin{Lemma}\label{l32}
Denote the left-hand side of ~\eqref{31} by $F$, i.e., $F:=\det\left(1-f K\right)_{\ell^2(\Z)}$. The following statements hold.
\begin{itemize}\itemsep=0pt
\item[{\rm (i)}]
$F_\ell$ does not depend on $\ell=0,1,2,\dots$ and is equal to $F$, i.e., we have $F_{\ell}=F$, for any~${\ell\in\Z_{\ge 0}}$.
\item[{\rm (ii)}] We have $\lim_{\ell \rightarrow\i} F_{\ell}=F_{\i}$.
\item[{\rm (iii)}] We have $K_{\i}(m_1,m_2)=-L(m_1,m_2)$
and thus $F_{\i}=\det(1+f L)_{\ell^2(\Z)}$.
\end{itemize}
\end{Lemma}
The proofs of the above (i)--(iii) will
be given in Sections \ref{s32}--\ref{s34}, respectively.

\begin{proof}[Proof of Theorem~\ref{t31}]
 This is a straightforward consequence of Lemma~\ref{l32}, which implies
 \begin{align*}
 \det(1-fK)_{\ell^2(\Z)}
 =
 F_{\ell}
 =
 F_{\i}
 =
 \det(1+fL)_{\ell^2(\Z)},
 \end{align*}
 under the assumptions given at the beginning of the section.
 \end{proof}

 Notice that one can significantly relax conditions on parameters set at the beginning of the section for Theorem~\ref{t31}. For this, one could relate matching of determinants to identity \eqref{t12-1}. Since the latter is an equality between polynomials of bounded degree in variables $a_1,\dots,a_N$, $b_1,\dots,b_N$ whose coefficients are rational functions in $q$, it can be extended to all choices of~$a_i$,~$b_j$ and in particular to all choices for which determinants in \eqref{31} make sense.


\subsection[Proof of Lemma 3.4 (i)]{Proof of Lemma~\ref{l32}\,(i)}
\label{s32}
 The case of $\ell=0$, i.e., $F_0=F$, was already mentioned after~\eqref{d31-2}.
 It is then sufficient to show~${F_{\ell}=F_{\ell-1}}$ for any $\ell=1,2,\dots$.
 For this purpose, we rewrite $F_{\ell}$ \eqref{d31-1} as a determinant with rank $N(\ell+ 1)$.
 \begin{Lemma}
 \label{l34}
 We have
 \begin{align}
 \label{3-12}
 F_{\ell}=\det\left(W_{n,n'}\right)_{n,n'=1,2,\dots,N(\ell+1)}.
 \end{align}
 Here $W_{n,n'}$ is defined by
 \begin{align}
\label{3-17}
 W_{n,n'}
 =
 \underset{w=\tilde{a}_{n}^{-1}}{\mathrm{Res}}
 \prod_{i=1}^N\frac{(b_i/w;q)_{\i}}{(a_i w;q)_{\i}}
 \int_{C_{n}}\frac{{\rm d}z}{z}
 \frac{-1}{z-\tilde{a}_{n'}^{-1}}
 \prod_{i=1}^N\frac{(a_i z;q)_{\i}}{(b_i/z;q)_{\i}}
 \frac{\theta\bigl(-\tilde{a}_{n}z/tq^{\frac12+k}\bigr)(q;q)_{\i}^2}{\theta\bigl(-1/tq^{\frac12+k}\bigr)\theta(\tilde{a}_n z)},
\end{align}
 where $C_n$ is a positively oriented circle with the condition $\tilde{a}_n^{-1}<|z|<(q\tilde{a}_n)^{-1}$ and $\theta(x)$ is defined below \eqref{a2-6}.
 \end{Lemma}


 \begin{proof}
 Let us write \eqref{3-10} as $f K_{\ell} = A_{\ell} B_{\ell}$, so that
 \[
 F_{\ell}=\det(1-fK_{\ell})_{\ell^2(\Z)}=\det(1-A_{\ell}B_{\ell})_{\ell^2(\Z)}.
 \]
 By Proposition~\ref{p33}\,(i), $A_{\ell}$, $B_{\ell}$ are
 Hilbert--Schmidt operators and hence $B_{\ell}A_{\ell}$ is trace class.
 Thus we can use the basic property of the determinant \smash{$\det(1-A_{\ell}B_{\ell})_{\ell^2(\Z)}=\det(\delta_{i,j}-(B_{\ell}A_{\ell})_{i,j})_{i,j=1}^{N(\ell+1)}$}
 and rewrite $F_{\ell}$ as
 \begin{align}
 \label{3-17d}
 F_{\ell}=\det(1-A_{\ell}B_{\ell})_{\ell^2(\Z)}
 =\det(\delta_{n,n'}-(B_{\ell}A_{\ell})_{n,n'})_{n,n'=1}^{N(\ell+1)}.
 \end{align}
Here
\begin{align*}%\label{3-15}
 (B_{\ell}A_{\ell})_{n,n'}
 &=
 \underset{w=\tilde{a}_n^{-1}}{\text{Res}}
 \prod_{i=1}^N\frac{(b_i/w;q)_{\i}}{(a_i w;q)_{\i}}
 \int_C\frac{{\rm d}z}{z}
 \frac{1}{z-\tilde{a}_{n'}^{-1}}
 \prod_{i=1}^N\frac{(a_i z;q)_{\i}}{(b_i/z;q)_{\i}}
 \sum_{m\in\Z}\frac{f(m)}{(z\tilde{a}_{n})^m}
 \\
 &=
 \underset{w=\tilde{a}_{n}^{-1}}{\text{Res}}
 \prod_{i=1}^N\frac{(b_i/w;q)_{\i}}{(a_i w;q)_{\i}}
 \int_{C_{n}'}\frac{{\rm d}z}{z}
 \frac{1}{z-\tilde{a}_{n'}^{-1}}
 \prod_{i=1}^N\frac{(a_i z;q)_{\i}}{(b_i/z;q)_{\i}}
 \cdot
 \frac{\theta\bigl(-\tilde{a}_{n}z/tq^{\frac12+k}\bigr)(q;q)_{\i}^2}{\theta\bigl(-1/tq^{\frac12+k}\bigr)\theta(\tilde{a}_n z)},
\end{align*}
 where in the second equality we applied a version of the Ramanujan summation formula \eqref{a2-6} and $\theta(x)$ is defined below \eqref{a2-6}. Note that to apply
 \eqref{a2-6}, we extend the contour $C$ to $C'_n$,
 which represents the circle centered at the origin with radius $R\in(q/\tilde{a}_n,1/\tilde{a}_n)$.

 Now we show that the entry $\delta_{n,n'}-(BA)_{n,n'}$ is equal to $W_{n,n'}$ \eqref{3-17}. For this, we further extend the radius $R$ of the contour such that $R\in(1/\tilde{a}_n,1/q\tilde{a}_n)$. Recalling the notation \eqref{s32-1} and the definition of $\theta (x)$ introduced below \eqref{a2-6} and
 noting \smash{$
 \prod_{i=1}^N(a_iz;q)_{\infty}=\prod_{\ell=1}^{\infty}(1-\tilde{a}_{\ell}z)$},
 we see
 \begin{gather}
 \underset{z=\tilde{a}_{n}^{-1}}{\text{Res}}
 \frac{{\rm d}z}{z}
 \frac{1}{z-\tilde{a}_{n'}^{-1}}
 \prod_{i=1}^N\frac{(a_i z;q)_{\i}}{(b_i/z;q)_{\i}}
 \cdot
 \frac{\theta\bigl(-\tilde{a}_{n}z/tq^{\frac12+k}\bigr)(q;q)_{\i}^2}{\theta\bigl(-1/tq^{\frac12+k}\bigr)\theta(\tilde{a}_n z)}\nonumber \\
 \qquad=\underset{z=\tilde{a}_{n}^{-1}}{\text{Res}}
 \frac{-\tilde{a}_{n'}}{z}
 \frac{\prod_{\substack{\ell=1\\ \ell\neq n'}}^{\infty}(1-\tilde{a}_\ell z)}{\prod_{i=1}^N(b_i/z;q)_{\i}}
 \cdot
 \frac{\theta\bigl(-\tilde{a}_{n}z/tq^{\frac12+k}\bigr)(q;q)_{\i}^2}{-\tilde{a}_n\bigl(z-\tilde{a}_n^{-1}\bigr)(q\tilde{a}_nz;q)_\infty \bigl(q\tilde{a}_n^{-1}/z;q\bigr)_\infty\theta\bigl(-1/tq^{\frac12+k}\bigr)}
 \nonumber\\
 \qquad = \tilde{a}_n
 \frac{\prod_{\substack{\ell=1\\ \ell\neq n'}}^{\infty}\bigl(1-\tilde{a}_\ell \tilde{a}_n^{-1}\bigr)}{\prod_{i=1}^N(b_i \tilde{a}_n;q)_{\i}}\delta_{n,n'}. \label{eq:zres}
 \end{gather}
 Here in the second equality we used the fact that in the second expression all factors do not include the residue at $z=\tilde{a}_n^{-1}$ except the factor $1/\big(z-\tilde{a}_n^{-1}\big)$ and this is eliminated by the factor~\smash{${\prod_{\substack{\ell=1\\ \ell\neq n'}}^{\infty}(1-\tilde{a}_\ell z)}$} unless $n=n'$.
 Similarly, we compute
 \begin{align}\label{eq:wres}
 \underset{w=\tilde{a}_{n}^{-1}}{\text{Res}}
 \prod_{i=1}^N\frac{(b_i/w;q)_{\i}}{(a_i w;q)_{\i}}=\underset{w=\tilde{a}_{n}^{-1}}{\text{Res}}\frac{\prod_{i=1}^N(b_i/w;q)_{\i}}{\prod_{\ell=1}^{\infty}(1-\tilde{a}_{\ell}w)}=-\tilde{a}_n^{-1}
 \frac{\prod_{i=1}^N(b_i \tilde{a}_n;q)_{\i}}{\prod_{\substack{\ell=1\\ \ell\neq n}}^{\infty}\big(1-\tilde{a}_\ell \tilde{a}_n^{-1}\big)}.
 \end{align}
 Therefore from \eqref{eq:zres} and \eqref{eq:wres}, the extension of the contour cancels the term $\delta_{n,n'}$ in the last expression in \eqref{3-17d} and we arrive at the expression \eqref{3-17}.
 \end{proof}


 Now we prove Lemma~\ref{l32}\,(i). We focus on the representation \eqref{3-12} and especially on the last ($N(\ell+1)$-th) row of the matrix $W$. We decompose the elements into two parts as
\begin{align*}%\label{3-18}
 W_{N(\ell+1),n'}
 &=W^{(1)}_{N(\ell+1),n'}+W^{(2)}_{N(\ell +1),n'}
\end{align*}
 for $n'\!=\!1,\dots,N(\ell+1)$. Here
 \smash{$W^{(1)}_{N(\ell+1),n'}$} corresponds to
 the residue at \smash{$z\!=\!\tilde{a}_{N(\ell+1)}^{-1}$}
 while~\smash{$W^{(2)}_{N(\ell+1),n'}$} corresponds to the element with the shrunk contour
 satisfying
 $q\tilde{a}_{n}^{-1}<|z|<\tilde{a}_{n}^{-1}$.
 As we stated in the last paragraph in the proof of Lemma~\ref{l34}, one easily sees \smash{$W^{(1)}_{N(\ell+1),n'}=\delta_{N(\ell+1),n'}$}.
 On the other hand recalling the notation \eqref{s32-1} and noting $q^{-1}\tilde{a}_n=\tilde{a}_{n-N}$, we have
 \begin{gather}
\underset{w=\tilde{a}_{n}^{-1}}{\text{Res}}
 \prod_{i=1}^N\frac{(b_i/w;q)_{\i}}{(a_i w;q)_{\i}}
 = q^{-1}
 \prod_{i=1}^N\frac{1}{1-b_i\tilde{a}_nq^{-1}}\frac{1}{1-a_i\tilde{a}_n^{-1}}
 \cdot
 \underset{w=\tilde{a}_{n-N}^{-1}}{\text{Res}}
 \prod_{i=1}^N\frac{(b_i/w;q)_{\i}}{(a_i w;q)_{\i}}, \label{eq:wn1}
 \\
 \theta(\tilde{a}_n x)=\frac{-q}{\tilde{a}_n x}
 \theta(\tilde{a}_{n-N}x).\label{eq:wn2}
 \end{gather}
 Applying \eqref{eq:wn1} to the first factor in \eqref{3-17} and \eqref{eq:wn2} to two factors $\theta\bigl(-\tilde{a}_n z/tq^{\frac12+k}\bigr)$ and $\theta(\tilde{a}_nz)^{-1}$ in the same equation, we see
 \[% \label{3-19}
 W^{(2)}_{N(\ell+1),n'}=-tq^{-\frac12+k}\prod_{i=1}^N\frac{1}{1-b_i\tilde{a}_{N(\ell+1)}q^{-1}}\frac{1}{1-a_i\tilde{a}_{N(\ell+1)}^{-1}}\cdot W_{N\ell,n'}.
 \]
 Combining them
 with the multi-linearity of determinants,
 we get
\[
 F_{\ell}=\det[W_{n,n'}]_{n,n'=1}^{N(\ell +1)}
 =
 \det[W_{n,n'}]_{n,n'=1}^{N(\ell+1)-1}.
\]
 Repeating this procedure $N-1$ times, we arrive at $F_{\ell}=F_{\ell-1}$.


\subsection[Proof of Lemma 3.4 (ii)]{Proof of Lemma~\ref{l32}\,(ii)}
\label{s33}
 From \eqref{3-10}
 and Corollary~\ref{cb2}, we find a uniform bound for the kernel
 \[
 \t(m_1)f(m_1)K_{\ell}(m_1,m_2)\t^{-1}(m_2),
 \]
 where
 $K_{\ell}(m_1,m_2)$, $\t(m)$ and $f(m)$ are defined by \eqref{d31-3}, \eqref{3-12d} and \eqref{p251-1} respectively. We have
 \begin{align}
 \label{3-32}
 \big|\t(m_1)f(m_1)K_{\ell}(m_1,m_2)\t^{-1}(m_2) \big|\le Dd^{|m_1|+|m_2|}
 \end{align}
 for some constants $D>0$ and $d\in (0,1)$ which do {\it not} depend on $\ell$, $m_1$ or $m_2$.

 Lemma~\ref{l32}\,(ii) is then immediately proven
 combining \eqref{3-32} with
 the Hadamard inequality
\begin{align}
\label{33-1}
 |\det A|\le
 \prod_{i=1}^n\Bigg(\sum_{j=1}^n|a_{i,j}|^2\Bigg)^{1/2},
\end{align}
which holds for any for $n\times n$ matrix $A=(a_{ij})_{i,j=1,\dots,n}$.
 Note that $F_\ell$ \eqref{d31-1} can be written as
\begin{align}
 F_{\ell}
\label{33-2}
 &=\sum_{n=0}^{\i}\frac{(-1)^n}{n!}
 \sum_{m_1=-\i}^{\i}\cdots\sum_{m_n=-\i}^{\i}
 \det\big(\t(m_i)\t^{-1}(m_j)f(m_i)K_\ell (m_i,m_j)\big)_{i,j=1}^n.
\end{align}
 By \eqref{3-32} and \eqref{33-1},
 the determinant in \eqref{33-2} can be bounded as
\begin{gather}
 \big|\det\big(\t(m_i)\t^{-1}(m_j)f(m_i)K_\ell (m_i,m_j)\big)_{i,j=1}^n \big|
 \nonumber\\
 \qquad \le
 \prod_{i=1}^n\Bigg(\sum_{j=1}^n \big|\t(m_i)\t^{-1}(m_j)f(m_i)K_\ell (m_i,m_j)\big|^2
 \Bigg)^{1/2}
 \nonumber\\
 \qquad\le
 \prod_{i=1}^n\Bigg(\sum_{j=1}^n \big|Dd^{|m_i|+|m_j|}\big|^2
 \Bigg)^{1/2}
 =D^n \prod_{i=1}^n d^{|m_i|}\Bigg(\sum_{j=1}^n d^{2|m_j|}\Bigg)^{1/2}\nonumber\\
 \qquad \le D^n n^{\frac{n}{2}}\prod_{i=1}^n d^{|m_i|}.\label{33-3}
\end{gather}
 Note that bound \eqref{33-3} holds for any $\ell\in\Z_{\ge 0}$. Moreover, we have
\begin{gather*}
 \sum_{m_1=-\i}^{\i}\cdots\sum_{m_n=-\i}^{\i}
 D^n n^{\frac{n}{2}}\prod_{i=1}^n d^{|m_i|}
 =D'^n n^{\frac{n}{2}} <\i,
 \qquad \sum_{n=0}^{\i}\frac{1}{n!}D'^n n^{\frac{n}{2}}<\i
\end{gather*}
with some constant $D'>0$.
 Thus by dominated convergence theorem and Proposition~\ref{p33}\,(ii), Lemma \ref{l32}
 (ii) holds.
\qed



\subsection[Proof of Lemma 3.4 (iii)]{Proof of Lemma~\ref{l32}\,(iii)}
\label{s34}

 In this subsection, we will
 prove $K_{\i}(m_1,m_2)
 =-L(m_1,m_2)$ by deforming the contour of $w$
 from $D_{\infty}$ to $\tilde{C}$ going through
 $\Gamma_c$, where all contours are depicted in Figure~\ref{f3}. We have to care about the behavior of the integrand $g_{a,b}(z,w;m_1,m_2)$ \eqref{p223}
 when $|w|$ is large since the poles~${w=1/a_iq^j}$,~${i=1,\dots,N}$, $j\in\Z_{\ge 0}$ accumulate at the infinite point of $w$ producing an essential singularity.
 We estimate the behavior for large $|w|$ using Appendix~\ref{ac}.


 Here for convenience we also use
 the notation \eqref{3-26} where
 we set $\sharp$ to be some contours
 depicted in
 Figure~\ref{f3}. Note that by definition, one sees
 $K_{\infty}(m_1,m_2)= K(m_1,m_2;D_{\i})$,
 and ${L(m_1,m_2)=-K\big(m_1,m_2;\tilde{C}\big)}$.
 Note that $\tilde{C}$ goes around the origin in negative direction.
 We now show that,
 from the properties of the $q$-Pochhammer symbol discussed in
 Appendix \ref{ac}, the following holds,
\begin{align}
\label{s34-2}
 K(m_1,m_2;\sharp)=0,
 \qquad \text{when}\quad \sharp=A_d,A_{-d},C_d,C_{-d}\ \text{and}\ \Gamma_{c_-}.
\end{align}
 In Figure~\ref{f3}, these contours
 are displayed in grey.
 These facts can be
 seen in the following way.
 In \eqref{3-26}, we focus on the
 integration of $w$,
 \begin{align}
 \label{3-29}
 \int_{\sharp}{\rm d}w\frac{w^{m}}{z-w}
 \prod_{i=1}^N
 \frac{{(b_i/w;q)_{\i}}}{(a_iw;q)_{\i}},
 \end{align}
 with $m\in\Z$ and $z\in C$ fixed. Here the contour $C$ is illustrated in Figure~\ref{f1}.
 Hereafter, we will explain the case $\sharp=\Gamma_{c_-}$ in detail whereas in other cases we will give an outline only in the paragraph below \eqref{3-34} since
 the other cases follow in a similar
 way.

 First we note that \eqref{3-29} does not depend on the value of $c_-$, as can be seen as follows. We consider the contour integral \eqref{3-29}
 with $\sharp=R_1\cup R_2\cup R_3\cup R_4$ where $R_i$, $i=1,2,3,4$ are depicted in
 Figure~\ref{f3}. Since there are no singularities inside the rectangle, the value of this integral is zero by the Cauchy's theorem. On the other hand
 from Lemma~\ref{lc1}\,(ii) we find the contributions from both $R_2$ and $R_4$ are bounded by $cd^m \exp\bigl[-Nc'\log^2 d\bigr]$ when $d>1$ for some constants $c$ and~$c'$, thus they vanish when $d\rightarrow\i$. Combining these facts, we find
 the contribution from $R_1$ is equal to that from $R_3$
 as $d\rightarrow\i$ up to sign.

 Next, we will show that \eqref{3-29} with $\sharp=\Gamma_{c_-}$ vanishes when $c_-\rightarrow -\i$. Combining Lem\-ma~\ref{lc1}\,(i) and (ii), we see that
 for $z=a+b\sqrt{-1}$ with $a<-1$ and $b\in\R$,
 \begin{align}
 \label{3-33}
 |(z;q)_{\i}|\ge c_1\exp\bigl[c_2\log^2(|a|\vee |b|)\bigr],
 \end{align}
 where $x\vee y=\max(x,y)$. Applying \eqref{3-33} to
 \eqref{3-29} with $\sharp=\Gamma_{c_-}$, we have
 \begin{align}
 \bigg| \int_{\Gamma_{c_-}} {\rm d}w h(w)\bigg|
 \le{}& \int_{\Gamma^{(1)}_{c_-}}{\rm d}w |h(w)|+\int_{\Gamma^{(2)}_{c_-}}{\rm d}w |h(w)|
 \nonumber\\
\le{}& c_3 |c_-|^{m_2}\exp\bigl(-Nc_2\log^2(|c_-|)\bigr) \nonumber\\
 &+c_4\int_{\Gamma_{c-}^{(2)}} {\rm d}w |w|^{m_2}
 \exp\bigl(-Nc_2\log^2|w|\bigr),
 \label{3-34}
 \end{align}
 where $h(w)$ is the integrand of \eqref{3-29}, $\Gamma_{c_-}^{(1)}$
 is the vertical line from $c_-+c_-\sqrt{-1}$ to $c_--c_-\sqrt{-1}$ and \smash{$\Gamma_{c_-}^{(2)}$} is its complement \smash{$\Gamma_{c_-}\setminus\Gamma_{c_-}^{(1)}$}.
 Thus we find that in the last expression in \eqref{3-34}, both two terms vanishes as $|c_-|\rightarrow\i$ and we arrive at the conclusion.

 We can also obtain \eqref{s34-2} with the other contours $\sharp=A_d,A_{-d},C_d,C_{-d}$ in a similar way, i.e., they follow from the following two facts: (I) The integration value \eqref{3-29} does not depend on the value of $d>1$ for any $m\in\mathbb{Z}$ and $z\in C$. (II) \eqref{3-29} vanishes as $d\rightarrow\infty$ for any~${m\in\mathbb{Z}}$ and $z\in C$. For proving (I), as in the case of $\Gamma_{c-}$, we introduce a positively-oriented rectangular contour~${\tilde{R}:=\tilde{R}_1\cup\tilde{R}_2\cup\tilde{R}_3\cup\tilde{R}_4}$, where $\tilde{R}_1$, $\tilde{R}_2$, $\tilde{R}_3$ and $\tilde{R}_4$ represents the right, top, left, and bottom edge respectively. In the case $A_d$ (resp.\ $A_{-d}$), we put $\tilde{R}$ such that its lower-left corner match the corner of $A_d$ (resp.\ $A_{-d}$). Then by using \eqref{lc1-3} in Lemma~\ref{lc1}\,(iii), we see that \eqref{3-29} with $\sharp=\tilde{R}_1$ vanishes as the length of the horizontal sides goes to infinity. In the case $C_d$ (resp.~$C_{-d}$), we set the length of the horizontal sides in $\tilde{R}$ to be $c-c_-$ and place it in a way that its two bottom (resp.\ top) corners match the corners in $C_d$ (resp.\ $C_{-d}$) and show that~\eqref{3-29} with~${\sharp=\tilde{R}_2}$ (resp.\ $\sharp{R}_4$) vanishes as the length of the vertical sides goes to infinity using Lemma~\ref{lc1} (ii). On the other hand one can show the property (II) just by using Lemma~\ref{lc1}\,(ii) and \eqref{lc1-3} in Lemma~\ref{lc1}\,(iii) in the cases $A_d$ and $A_{-d}$ while in the case $C_{d}$ and~$C_{-d}$, we use Lemma~\ref{lc1}\,(ii).


 Now we will prove Lemma \ref{l32}\,(iii)
 by showing the following two relations:
\begin{subequations}
 \begin{align}
\label{s34-5}
 &K(m_1,m_2;\Gamma_c)
 =K_{\i}(m_1,m_2),
 \\
\label{s34-6}
 &K(m_1,m_2;\Gamma_c)
 =-L(m_1,m_2).
 \end{align}
\end{subequations}
 First, we prove \eqref{s34-5}. Noting $\Gamma_c=A_d\cup D_{\infty}\cup A_{-d}$, and \eqref{s34-2}, we have
\begin{align*}%\label{s34-4}
 K(m_1,m_2;\Gamma_c)
 &=K(m_1,m_2;D_{\infty})+K(m_1,m_2;A_d)+K(m_1,m_2;A_{-d})\\
 &=K(m_1,m_2;D_{\infty}).
\end{align*}
 Recalling by definition $K(m_1,m_2; D_{\infty})=K_{\infty}(m_1,m_2)$, we arrive at \eqref{s34-5}.

 Next we show the relation \eqref{s34-6}.
 Noting the decomposition
 $\tilde{C}=\Gamma_{c}\cup C_d\cup C_{-d}\cup\Gamma_{c_-}$, and~\eqref{s34-2}, we have
 \[
 K\big(m_1,m_2,\tilde{C}\big)=
 K(m_1,m_2,\Gamma_c).
 \]
 Combining this with the fact $L(m_1,m_2)=-K\big(m_1,m_2;\tilde{C}\big)$, we arrive at \eqref{s34-6}.
\qed

\begin{figure}[ht]
 \centering
 \includegraphics[scale=1]{fig3.pdf}
 %\includegraphics[]{contours_IMS21_1.pdf}
 \caption{All of the contours appearing in the proof of
 Lemma \ref{l32}. We set
 $c>0$, $c_-<0$. }
 \label{f3}
\end{figure}


